%-------------------------
% Resume in Latex (Chinese Version)
%------------------------

\documentclass[letterpaper,11pt]{article}

\usepackage{latexsym}
\usepackage[empty]{fullpage}
\usepackage{titlesec}
\usepackage{marvosym}
\usepackage[usenames,dvipsnames]{color}
\usepackage{verbatim}
\usepackage{enumitem}
\usepackage{hyperref}
\usepackage{fancyhdr}
\usepackage{xeCJK}
\setCJKmainfont{Noto Serif CJK TC}
\setCJKsansfont{Noto Sans CJK TC}
\setCJKmonofont{Noto Sans Mono CJK TC}


\pagestyle{fancy}
\fancyhf{} % clear all header and footer fields
\fancyfoot{}
\renewcommand{\headrulewidth}{0pt}
\renewcommand{\footrulewidth}{0pt}

% Adjust margins
\addtolength{\oddsidemargin}{-0.375in}
\addtolength{\evensidemargin}{-0.375in}
\addtolength{\textwidth}{1in}
\addtolength{\topmargin}{-.5in}
\addtolength{\textheight}{1.0in}

\urlstyle{same}

\raggedbottom
\raggedright
\setlength{\tabcolsep}{0in}

% Sections formatting
\titleformat{\section}{
  \vspace{-4pt}\scshape\raggedright\large
}{}{0em}{}[\color{black}\titlerule \vspace{-5pt}]

%-------------------------
% Custom commands
\newcommand{\resumeItem}[2]{
  \item\small{
    \textbf{#1}{: #2 \vspace{-2pt}}
  }
}

\newcommand{\resumeSubheading}[4]{
  \vspace{-1pt}\item
    \begin{tabular*}{0.97\textwidth}{l@{\extracolsep{\fill}}r}
      \textbf{#1} & #2 \\
      \textit{\small#3} & \textit{\small #4} \\
    \end{tabular*}\vspace{-5pt}
}

\newcommand{\resumeSubItem}[2]{\resumeItem{#1}{#2}\vspace{-4pt}}

\renewcommand{\labelitemii}{$\circ$}

\newcommand{\resumeSubHeadingListStart}{\begin{itemize}[leftmargin=*]}
\newcommand{\resumeSubHeadingListEnd}{\end{itemize}}
\newcommand{\resumeItemListStart}{\begin{itemize}}
\newcommand{\resumeItemListEnd}{\end{itemize}\vspace{-5pt}}

%-------------------------------------------
%%%%%%  CV STARTS HERE  %%%%%%%%%%%%%%%%%%%%%%%%%%%%


\begin{document}

%----------HEADING-----------------
\begin{tabular*}{\textwidth}{l@{\extracolsep{\fill}}r}
  \textbf{\href{https://github.com/twohorse0311}{\Large 馮元杰（Francis）}} & 信箱：\href{mailto:fong54388@gmail.com}{fong54388@gmail.com}\\
  \href{https://www.linkedin.com/in/yuan-jie-fong-9289061b4}{https://www.linkedin.com/in/yuan-jie-fong-9289061b4} & 電話：0921803781 \\
  全端工程師 & \\
\end{tabular*}


%-----------EXPERIENCE-----------------
\section{工作經歷}
  \resumeSubHeadingListStart

    \resumeSubheading
      {薩泰爾娛樂（STR Network）}{台北，台灣}
      {軟體工程師}{2024 年 1 月 -- 至今}
      \resumeItemListStart
        \resumeItem{\href{https://inviti.vip/}{貴賓管理系統}}
          {使用 Django、PostgreSQL、WebSocket 與 Redis 設計並建置即時貴賓管理系統，部署於 GCP（透過 Zeabur）。運用 DDD 方法論與跨部門利害關係人協作，確保系統貼合實際業務流程。實現 10--50 名內部使用者的即時資料同步，降低禮賓統籌超過 90\% 的工作量。}
        \resumeItem{內部會計系統}
          {使用 Django 與 PostgreSQL 開發涵蓋核銷、匯款及專案預付款請款結帳的全流程會計系統。深入營運與財務團隊進行使用者研究以釐清實際工作流程。降低超過 70\% 的處理時間，同時達成現金追蹤零誤差的準確度。}
        \resumeItem{\href{https://www.businessweekly.com.tw/event/ai100/}{商業週刊 AI 創新百強}}
          {主導公司 AI 相關開發專案，協助薩泰爾娛樂入選商業週刊 AI 創新百強企業。}
      \resumeItemListEnd

    \resumeSubheading
      {國立臺北大學統計學系}{台北，台灣}
      {課程講師}{2025 年 9 月 -- 2025 年 12 月}
      \resumeItemListStart
        \resumeItem{Python 全端開發課程}
          {為大學部學生設計並教授以 Python 與 Django 為核心的全端網頁開發課程。}
      \resumeItemListEnd

    \resumeSubheading
      {\href{https://www.amiino.com.tw/}{美無痕生物科技}}{台北，台灣}
      {廣告數據分析師}{2025 年 3 月 -- 至今}
      \resumeItemListStart
        \resumeItem{廣告數據分析平台}
          {使用 Django 建置全端分析平台，整合 3 家廣告商共 8 個渠道（Google 關鍵字廣告/Pmax/YouTube、Facebook、Criteo、Yahoo、Bing）。建立 Data ETL 流程，將跨渠道資料彙整至 PostgreSQL，並搭配自動化視覺化儀表板進行趨勢分析與綜合比較。}
        \resumeItem{AI 自動化報告}
          {整合 AI 自動產生每週數據報告，比較跨期間績效表現，助力 2025 年度 ROAS 成長超過 50\%。}
      \resumeItemListEnd

    \resumeSubheading
      {個人工作室（影視產業）}{台北，台灣}
      {數據分析師}{2021 年 4 月 -- 至今}
      \resumeItemListStart
        \resumeItem{數據分析與視覺化}
          {建立自動化網路爬蟲從平台 API 擷取銷售數據，建置本地資料庫，並自動化資料清洗、視覺化與月報產出。}
        \resumeItem{成效}
          {加入後月營收成長 10 倍，2 年內累計營收超過 200,000 美元。}
      \resumeItemListEnd


  \resumeSubHeadingListEnd


%-----------EDUCATION-----------------
\section{學歷}
  \resumeSubHeadingListStart
    \resumeSubheading
      {國立清華大學 MBA}{新竹，台灣}
      {主修資訊系統與商業分析}{2022 年 9 月 -- 2024 年 8 月}
    \resumeSubheading
      {國立臺北大學}{台北，台灣}
      {統計學系 工商管理學士}{2017 年 9 月 -- 2021 年 6 月}
  \resumeSubHeadingListEnd


%-----------PROJECTS-----------------
\section{獎項與專案經歷}
  \resumeSubHeadingListStart
    \resumeSubItem{\href{https://webconf.tw/}{WebConf Taiwan 2025 講者}}{於 WebConf Taiwan 2025 發表「AI 時間槓桿賦能之術——全端工程師獨自升級全記錄」，分享運用 AI 提升工程生產力的實戰經驗。}
  \resumeSubHeadingListEnd

%
%--------PROGRAMMING SKILLS------------
\section{技能摘要}
 \resumeSubHeadingListStart
   \item{
     \textbf{程式語言}{：Python, SQL, Ruby, JavaScript, R}
   }
   \item{
     \textbf{框架與工具}{：Django, Flask, Redis, WebSocket, PostgreSQL, Git, LaTeX}
   }
   \item{
     \textbf{雲端與基礎設施}{：GCP, Zeabur}
   }
   \item{
     \textbf{語言能力}{：中文母語、英文流利}
   }
 \resumeSubHeadingListEnd


%-------------------------------------------
\end{document}
